% p210_1; source PDF page 210; manually verified editable table.
\begin{tikzpicture}[x=1cm,y=-1cm,inner sep=0pt,outer sep=0pt]
\path[use as bounding box] (0,0) rectangle (16.000,4.600);
\draw[line width=.4pt] (0.000,0.000) rectangle (7.000,1.600);
\node[align=left,text width=6.740cm,font={\linespread{1}\fontsize{10}{13.50}\selectfont}] at (3.500,0.800) {缓刑考验期内发现漏罪，假释考验期内发现漏罪，数罪并罚后};
\draw[line width=.4pt] (7.000,0.000) rectangle (16.000,1.600);
\node[align=left,text width=8.740cm,font={\linespread{1}\fontsize{10}{13.50}\selectfont}] at (11.500,0.800) {可以再次缓刑，可以再次假释；因为发现漏罪不表示犯罪人的人身危险性升高};
\draw[line width=.4pt] (0.000,1.600) rectangle (7.000,2.600);
\node[align=left,text width=6.740cm,font={\linespread{1}\fontsize{10}{13.50}\selectfont}] at (3.500,2.100) {缓刑考验期内犯新罪，数罪并罚后};
\draw[line width=.4pt] (7.000,1.600) rectangle (16.000,2.600);
\node[align=left,text width=8.740cm,font={\linespread{1}\fontsize{10}{13.50}\selectfont}] at (11.500,2.100) {不能再次缓刑};
\draw[line width=.4pt] (0.000,2.600) rectangle (7.000,3.600);
\node[align=left,text width=6.740cm,font={\linespread{1}\fontsize{10}{13.50}\selectfont}] at (3.500,3.100) {假释考验期内犯新罪，数罪并罚后};
\draw[line width=.4pt] (7.000,2.600) rectangle (16.000,3.600);
\node[align=left,text width=8.740cm,font={\linespread{1}\fontsize{10}{13.50}\selectfont}] at (11.500,3.100) {不能再次假释};
\draw[line width=.4pt] (0.000,3.600) rectangle (7.000,4.600);
\node[align=left,text width=6.740cm,font={\linespread{1}\fontsize{10}{13.50}\selectfont}] at (3.500,4.100) {缓刑考验期内犯新罪，数罪并罚后};
\draw[line width=.4pt] (7.000,3.600) rectangle (16.000,4.600);
\node[align=left,text width=8.740cm,font={\linespread{1}\fontsize{10}{13.50}\selectfont}] at (11.500,4.100) {可以假释};
\end{tikzpicture}
